\documentclass[11pt]{article}
\usepackage[T1]{fontenc}
\usepackage{lmodern}
\usepackage[margin=1in]{geometry}
\usepackage{amsmath,amssymb,amsthm,booktabs,array,longtable}
\usepackage[colorlinks=true,linkcolor=blue,urlcolor=blue,citecolor=blue]{hyperref}
\newtheorem{theorem}{Theorem}
\newtheorem{corollary}[theorem]{Corollary}
\newcommand{\mob}{\mu}
\newcommand{\rad}{\operatorname{rad}}
\newcommand{\C}{C}
\title{Convention corrections and alternative proofs for\\
Morken's conjectures in OEIS A112399\\
\large Revision 41, FORMULA; after Sander--Sander (2009)\\
and Ramar\'e (2014--2015)}
\author{Haobo Ma\\\small ChronoAI Pte Ltd, Singapore
\and Wenlin Zhang\\\small National University of Singapore\\
\small ChronoAI Pte Ltd, Singapore}
\date{4 November 2026}
\begin{document}
\maketitle
\begin{abstract}
We give alternative proofs of Jamie Morken's two conjectures in OEIS
A112399, revision 41, FORMULA (18 September 2026), following the finite
coprimality identity of Sander and Sander (2009) and equivalent identities
of Ramar\'e (2014--2015) \cite{a112399,sander,ramare2014,ramare2015}.
Both coprime M\"obius sum identities hold with an arbitrary nonnegative
integer cutoff. Complete finite derivations cover the A400203 triangle,
the rectangular A372700 identity, and nearby divisor and gcd transforms,
with their earlier proofs and formula records distinguished. We correct
the strict endpoint in A304575's name and the signed-variant convention
in A112526. The product identity is reproved with credit to Bala (2026).
\end{abstract}

\section*{Earlier proofs and scope}
The first identity below is a specialization of the finite coprimality
expansion in Sander and Sander \cite[proof of Proposition 4, (5)--(7)]{sander}:
choose the vector entries to be M\"obius values up to the cutoff and zero
after it, in an ambient dimension at least the parameter and the cutoff.
Ramar\'e \cite[\S1, immediately after (1.2)]{ramare2014} gives an equivalent
reciprocal-weight identity. Taking its jump at each positive integer $r$,
multiplying by $r$, and summing up to the cutoff gives \eqref{eq:first}.
For the second identity, Ramar\'e's \cite[\S16, (16.1)]{ramare2015}
squarefree-divisor and prime-supported-square decomposition becomes
\eqref{eq:second} under the unique factorization $q=dw^2$, with $d$
squarefree. That paper was electronically published on 1 December 2014.

Ingix and reuns \cite{mse2019} give a related coefficient proof on
13 June 2019; reuns's corrected support condition is $\rad(m)\mid a$,
not $m\mid a$. Helfgott \cite[\S5.3.4, (5.73)]{helfgott} proves the exact
coefficient cancellation and a weighted summation identity. The same coefficient and
weighted identities appear in Ramar\'e--Zuniga-Alterman
\cite[Lemma 2.5]{rz2026}, first posted on 26 March 2026, and in
Zuniga-Alterman \cite[Lemma 7.1]{z2026}, posted on 13 October and revised
on 26 October 2026, after the posing date. In each weighted identity,
choosing $h(t)=t\mathbf1_{t\le X}$ makes the specialization finite.

The arbitrary-function transform at the diagonal has an earlier
reduced-fraction proof by West \cite{west}; the finite gcd formula is
a specialization of Schramm \cite[theorem (i), (2), and Example 5]{schramm}.
The rectangular identity follows from the same finite coprimality
expansion. A400203's formulas and the three triangle specializations
below were recorded by Morken on 21 September 2026 in its history,
revisions 2 and 4 \cite{a400203}; 1 October dates the later revision-17
snapshot, not their first entry. The two corrected source conventions
are stated precisely in Corollaries~\ref{cor:signed} and \ref{cor:counts}.

\section{The two identities}
Write $\mob$ for the integer-valued M\"obius function, with $\mob(1)=1$;
thus $\mob(r)=0$ when a prime square divides $r$, and otherwise
$\mob(r)=(-1)^s$ when $r$ has $s$ distinct prime factors \cite{a008683}.
For integers $n\geq1$ and $X\geq0$, define
\[
 \C_n(X)=\sum_{\substack{1\leq k\leq X\\(k,n)=1}}\mob(k),
 \qquad M(X)=\sum_{1\leq k\leq X}\mob(k),
 \qquad \rad(n)=\prod_{p\mid n}p.
\]
Here $(k,n)$ denotes the greatest common divisor, $\rad(1)=1$, and empty
sums are zero. The functions $M$ and $\rad$ are OEIS A002321 and A007947
\cite{a002321,a007947}. The sequence A112399 is $a(n)=\C_n(n)$.
Its two conjectures assert the $X=n$ cases of the following theorem.

\begin{theorem}\label{thm:main}
For every integer $n\geq1$ and every integer $X\geq0$,
\begin{align}
 \C_n(X)
 &=\sum_{\substack{d\mid n\\\mob(d)\ne0}}
       \C_d(\lfloor X/d\rfloor),\label{eq:first}\\
 \C_n(X)
 &=\sum_{\substack{1\leq q\leq X\\\rad(q)\mid n}}
       M(\lfloor X/q\rfloor).\label{eq:second}
\end{align}
All divisors and summation indices are positive.
\end{theorem}

The condition $\mob(d)\ne0$ restricts the outer divisor, not $n$ itself.
In \eqref{eq:first}, the inner sum uses $(m,d)=1$, whereas the left side
uses $(k,n)=1$. In \eqref{eq:second}, $\rad(q)\mid n$ allows arbitrary
powers of primes dividing $n$; it does not require $q\mid n$.

\begin{proof}
We first record the pointwise identity, for positive $d,m$,
\begin{equation}\label{eq:pointwise}
 \mob(d)\mob(dm)=
 \begin{cases}
   \mob(m),&\mob(d)\ne0\text{ and }(m,d)=1,\\
   0,&\text{otherwise}.
 \end{cases}
\end{equation}
If $\mob(d)=0$, the left side vanishes. Otherwise $d$ is squarefree.
When $(m,d)=1$, multiplicativity gives
$\mob(dm)=\mob(d)\mob(m)$ and $\mob(d)^2=1$.
When $(m,d)>1$, a common prime has its square dividing $dm$, so
$\mob(dm)=0$. This proves \eqref{eq:pointwise} in all cases.

Apply \eqref{eq:pointwise}, put $k=dm$, and exchange finite sums:
\begin{align*}
 \sum_{\substack{d\mid n\\\mob(d)\ne0}}
       \C_d(\lfloor X/d\rfloor)
 &=\sum_{d\mid n}\mob(d)\sum_{1\leq m\leq X/d}\mob(dm)\\
 &=\sum_{1\leq k\leq X}\mob(k)
        \sum_{d\mid(k,n)}\mob(d).
\end{align*}
For any positive $g$, the inner divisor sum equals
$\prod_{p\mid g}(1-1)$: it is $1$ if $g=1$ and $0$ otherwise.
The last display therefore equals $\C_n(X)$, proving \eqref{eq:first}.
The reindexing uses $m\leq\lfloor X/d\rfloor$ if and only if $dm\leq X$,
so it also covers $d>X$ and $X=0$.

For \eqref{eq:second}, define $h_n(q)=1$ when $\rad(q)\mid n$ and
$h_n(q)=0$ otherwise. For every positive integer $r$ we claim
\begin{equation}\label{eq:convolution}
 \sum_{q\mid r}h_n(q)\mob(r/q)
   =\begin{cases}\mob(r),&(r,n)=1,\\0,&(r,n)>1.\end{cases}
\end{equation}
To check this, factor the divisor sum prime by prime. If $p^a\Vert r$
and $p\nmid n$, only exponent zero is allowed in $q$, so the local
factor is $\mob(p^a)$. If $p\mid n$, every exponent $0\leq j\leq a$
is allowed, and the local factor is
\[
 \sum_{j=0}^{a}\mob(p^{a-j})=1-1=0\qquad(a\geq1).
\]
For $r=1$ the empty product is $1$. These observations give
\eqref{eq:convolution}. Summing it for $1\leq r\leq X$ and putting
$r=qk$ yields
\[
 \C_n(X)=\sum_{1\leq r\leq X}\sum_{q\mid r}h_n(q)\mob(r/q)
 =\sum_{1\leq q\leq X}h_n(q)
       \sum_{1\leq k\leq\lfloor X/q\rfloor}\mob(k),
\]
which is \eqref{eq:second}.
\end{proof}

For example, $a(10)=-1$. In \eqref{eq:first}, the contributions from
$d=1,2,5,10$ are respectively $M(10)=-1$, $\C_2(5)=-1$,
$\C_5(2)=0$, and $\C_{10}(1)=1$. In \eqref{eq:second}, the allowed
$q\leq10$ are $1,2,4,5,8,10$, giving $-1,-2,0,0,1,1$.
For $n=1$, both formulas reduce to $M(X)$; for $X=0$, both are empty.

\section{Product and quadratic identities}
\begin{corollary}[Product formula; Bala's weighted identity]\label{cor:product}
For a multiplicative arithmetic function $f$, $n\geq1$ and $X\geq0$,
\[
 \sum_{k=1}^{X}\mob(nk)f(nk)
 =\mob(n)f(n)\sum_{\substack{k\leq X\\(k,n)=1}}\mob(k)f(k).
\]
In particular,
\begin{equation}\label{eq:product}
 \sum_{1\leq k\leq X}\mob(nk)=\mob(n)\C_n(X).
\end{equation}
\end{corollary}
\begin{proof}
For squarefree $n$ the terms with $(k,n)>1$ vanish and the others
factor by multiplicativity. For nonsquarefree $n$, both sides vanish.
\end{proof}
The $X=n$ instance is Morken's conjecture in A372706 \cite{a372706}.
Peter Bala proved it on 22 September 2026; his Proposition~1 also allows
a multiplicative weight $f(nk)$ \cite{bala}. Equation \eqref{eq:product}
is its arbitrary-cutoff specialization. Among the primary proofs located,
Bala's dated proof is the earliest for this product statement; the two
A112399 identities concern different divisor sums.

\begin{corollary}[Rectangular quadratic identity]\label{cor:quadratic}
For integers $X,Y\geq0$,
\begin{equation}\label{eq:rectangle}
 \sum_{j=1}^{X}\sum_{k=1}^{Y}\mob(jk)
 =\sum_{d=1}^{\min(X,Y)}\mob(d)
   \C_d(\lfloor X/d\rfloor)\C_d(\lfloor Y/d\rfloor).
\end{equation}
Thus, for $T(N)=\sum_{j,k=1}^{N}\mob(jk)$ (A372700 \cite{a372700}),
\begin{equation}\label{eq:quadratic}
 T(N)=\sum_{d=1}^{N}\mob(d)\C_d(\lfloor N/d\rfloor)^2
 \qquad(N\geq1).
\end{equation}
\end{corollary}
\begin{proof}
For positive $d,x,y$,
\[
 \mob(d)\mob(x)\mob(y)\mathbf1_{(x,d)=(y,d)=1}
       =\mob(d)\mob(dx)\mob(dy).
\]
If $d$ is nonsquarefree, both sides vanish; otherwise use
multiplicativity when both coprimality conditions hold, and squarefree
vanishing when either fails. Expanding the right side of
\eqref{eq:rectangle} and setting $j=dx$, $k=dy$ gives
\[
 \sum_{\substack{j\leq X\\k\leq Y}}\mob(j)\mob(k)
       \sum_{d\mid(j,k)}\mob(d)
   =\sum_{\substack{j\leq X,\ k\leq Y\\(j,k)=1}}\mob(j)\mob(k)
   =\sum_{\substack{j\leq X\\k\leq Y}}\mob(jk).
\]
The last equality again uses squarefree vanishing for noncoprime $j,k$.
Setting $X=Y=N$ proves Morken's quadratic formula.
\end{proof}

\begin{corollary}[Row sum and plateaux]\label{cor:row}
For $N\geq1$,
\begin{equation}\label{eq:row}
 T(N)=-1+2\sum_{k=1}^{N}\mob(k)a(k).
\end{equation}
In particular $T(N)=T(N-1)$ when $N$ is nonsquarefree.
\end{corollary}
\begin{proof}
$T(1)=1$. For $N>1$, the last row and column give
$T(N)-T(N-1)=2\sum_{k=1}^{N}\mob(Nk)-\mob(N^2)
=2\mob(N)a(N)$, by Corollary~\ref{cor:product} and $\mob(N^2)=0$.
Summation gives \eqref{eq:row}, including its initial value.
\end{proof}
The A372700 conjecture of infinitely many sign changes is not covered.

\begin{corollary}[All finite formulas of A400203]\label{cor:triangle}
Let $U(N,d)=\C_d(\lfloor N/d\rfloor)$ for $1\leq d\leq N$,
the triangle A400203, revision 17, FORMULA \cite{a400203}.
Then
\begin{equation}\label{eq:triangle-mertens}
 U(N,d)=\sum_{\substack{q\leq N/d\\\rad(q)\mid d}}
                  M(\lfloor N/(dq)\rfloor).
\end{equation}
Its boundary and plateau formulas are
\[
 U(N,1)=M(N),\quad U(N,d)=1\ (2d>N),\quad
 U(N,d)=U(d\lfloor N/d\rfloor,d).
\]
For $1\leq d<N$, its increment is $\mob(N/d)$ when
$d\mid N$ and $(N/d,d)=1$, and zero otherwise.
Writing $H(k)=1$ if every prime exponent of $k$ is at least two,
and $H(k)=0$ otherwise, with $H(1)=1$, gives
\begin{align}
 \sum_{d=1}^{N}U(N,d)&=\sum_{k=1}^{N}H(k)
                         =\mathrm{A217038}(N+1),\label{eq:triangle-row}\\
 \sum_{d=1}^{N}\mob(d)U(N,d)&=\sum_{k=1}^{N}\mob(k)\tau(k)
                         =\mathrm{A062563}(N),\label{eq:triangle-mu}\\
 \sum_{d=1}^{N}\mob(d)^2U(N,d)&=1.\label{eq:triangle-square}
\end{align}
Here $H$ is A112526, $\tau(k)$ is the number of positive divisors,
and A217038 counts powerful numbers strictly below its argument
\cite{triangle-aux}.
\end{corollary}
\begin{proof}
Equation \eqref{eq:triangle-mertens} is \eqref{eq:second} with parameter
$d$ and cutoff $\lfloor N/d\rfloor$, using
$\lfloor\lfloor N/d\rfloor/q\rfloor=\lfloor N/(dq)\rfloor$.
The boundary, plateau and increment formulas follow by considering
whether $\lfloor N/d\rfloor$ changes and whether its new endpoint is
coprime to $d$. In particular $U(N,N)=1$.

Reindex each row by $k=dx$. For \eqref{eq:triangle-row}, its coefficient
at $k$ is $\sum_{x\mid k,(x,k/x)=1}\mob(x)$. At a prime power $p^e$
the two choices $x=1,p^e$ give $1+\mob(p^e)$: zero for $e=1$ and
one for $e\geq2$. The product is $H(k)$.
For \eqref{eq:triangle-mu} the coefficient is
$\sum_{dx=k,(d,x)=1}\mob(d)\mob(x)$. It is zero if $k$ is
nonsquarefree, and otherwise every divisor contributes $\mob(k)$,
giving $\mob(k)\tau(k)$; equivalently it is $\mob(k)2^{\omega(k)}$.
For \eqref{eq:triangle-square}, the local factor of
$\sum_{dx=k,(d,x)=1}\mob(d)^2\mob(x)$ is
$\mob(p^e)+\mob(p^e)^2=0$ for every $e\geq1$.
Thus only $k=1$ contributes, with coefficient one.
\end{proof}
The prime-supported formula \eqref{eq:triangle-mertens} is already recorded
in A400203, revision 17 (1 October 2026); Theorem~\ref{thm:main} supplies
its cancellation proof. Taking $N=dX$ recovers any nonnegative cutoff.

\begin{corollary}[A112526's signed-variant convention]\label{cor:signed}
With the powerful-number indicator $H$ above,
$H(n)\mob(n)=\mathbf1_{n=1}$.
The multiplicative function with prime-power values zero for exponent
one and $-1$ for larger exponents is instead $H(n)\mob(\rad(n))$.
\end{corollary}
\begin{proof}
Every powerful $n>1$ is nonsquarefree, so $\mob(n)=0$.
For the second assertion, evaluate at $p^e$ and use multiplicativity.
This corrects the pointwise-product description of the signed variant
in A112526, COMMENT \cite{triangle-aux}.
\end{proof}

\section{Divisor grouping and further siblings}
\begin{corollary}[A general finite transform]\label{cor:group}
Let $F$ be any complex-valued function on positive integers. For $n\geq1$
and $X\geq0$,
\begin{equation}\label{eq:group}
 \sum_{k=1}^{X}F\!\left(\frac{k}{(k,n)}\right)
 =\sum_{d\mid n}\ \sum_{\substack{1\leq m\leq\lfloor Xd/n\rfloor\\(m,d)=1}}F(m).
\end{equation}
In particular, if $A_F(d)=\sum_{m\leq d,(m,d)=1}F(m)$ and
$B_F(n)=\sum_{k\leq n}F(k/(k,n))$, then
\begin{equation}\label{eq:general-transform}
 B_F(n)=\sum_{d\mid n}A_F(d),\qquad
 A_F(n)=\sum_{d\mid n}\mob(n/d)B_F(d).
\end{equation}
\end{corollary}
\begin{proof}
Group by $g=(k,n)$, put $d=n/g$ and $k=gm$.
Then $(m,d)=1$, and $k\leq X$ is equivalent to
$m\leq\lfloor Xd/n\rfloor$. This is a bijection in both directions,
proving \eqref{eq:group}. For $X=n$ it gives the first transform;
substitution and $\sum_{e\mid t}\mob(e)=\mathbf1_{t=1}$ give its inverse.
\end{proof}

\begin{corollary}[A332685 and the A112399 transform]\label{cor:transform}
For $b(n)=\sum_{k\leq n}\mob(k/(k,n))$ (A332685 \cite{a332685}),
\begin{equation}\label{eq:transform}
 b(n)=\sum_{d\mid n}a(d),\qquad
 a(n)=\sum_{d\mid n}\mob(n/d)b(d).
\end{equation}
Also $b(n)=\sum_{k\leq n}\mob(\operatorname{lcm}(n,k)/n)$.
\end{corollary}
\begin{proof}
Take $F=\mob$ in Corollary~\ref{cor:group}; use
$\operatorname{lcm}(n,k)/n=k/(k,n)$ for the last expression.
\end{proof}
The inverse formula is credited to Amiram Eldar in A112399;
the forward formula and lcm expression are recorded in A332685.

\begin{corollary}[Squarefree counts and the A304575 endpoint]\label{cor:counts}
Write $s(d)=\sum_{k\leq d,(k,d)=1}\mob(k)^2$ (A073311 \cite{a073311}).
The formulas and data of A304575 \cite{a304575} give the sequence
\begin{equation}\label{eq:counts}
 R(n)=\sum_{d\mid n}s(d)
 =\sum_{k=1}^{n}\mob\!\left(\frac{k}{(k,n)}\right)^2
 =\sum_{k=1}^{n}\mob\!\left(\frac{\operatorname{lcm}(n,k)}{n}\right)^2.
\end{equation}
Consequently $R(n)\leq n$ and
$s(n)=\sum_{d\mid n}\mob(n/d)R(d)$.
If the strict inequality $k<d$ in that entry's NAME is read literally
for positive $k$, its divisor sum equals $R(n)-1$, not $R(n)$.
For $Q(t)=\sum_{1\leq k<t}\mob(k)^2$ (A013928 \cite{a013928}),
every prime $p$ and integer $m\geq0$ satisfy
\begin{align}
 R(p^{2m})&=1+\sum_{i=0}^{m}Q(p^{2i}),\label{eq:even}\\
 R(p^{2m+1})&=1+\sum_{i=0}^{m}Q(p^{2i+1}).\label{eq:odd}
\end{align}
\end{corollary}
\begin{proof}
Apply Corollary~\ref{cor:group} with $F=\mob^2$, and use the lcm
identity. Each summand in the middle expression of \eqref{eq:counts}
is zero or one, so $R(n)\leq n$.
For $d>1$ the endpoint $k=d$ is not coprime to $d$; for $d=1$ it
contributes one. Every divisor sum therefore differs by exactly one
under the strict convention. In particular the literal NAME gives zero
at $n=1$, whereas the entry's data and formulas give $R(1)=1$.

Put $D_j=s(p^j)$, with $D_0=1$ and $D_1=Q(p)$.
For $j\geq2$, a squarefree integer below $p^j$ either avoids $p$
(there are $D_j$ such integers) or equals $p\ell$, where $\ell<p^{j-1}$
is squarefree and avoids $p$ (there are $D_{j-1}$).
Thus $D_j+D_{j-1}=Q(p^j)$. Pair consecutive terms in
$R(p^e)=\sum_{j=0}^{e}D_j$; the unpaired terms are $D_0=1$
for even $e$, and $D_0+D_1=1+Q(p)$ for odd $e$.
This gives \eqref{eq:even}--\eqref{eq:odd}, also at $m=0$ since $Q(1)=0$.
\end{proof}
The inclusive convention in \eqref{eq:counts} agrees with A073311 and
with A304575's formulas and programs. Only the strict-endpoint wording
of its NAME requires correction.

\begin{corollary}[Other finite gcd sums]\label{cor:gcd}
For $n\geq1$ and any arithmetic function $F$,
\begin{align}
 \sum_{k=1}^{n}F((k,n))&=\sum_{d\mid n}F(d)\varphi(n/d),\label{eq:gcd}\\
 \sum_{k=1}^{n}F(n/(k,n))&=\sum_{d\mid n}F(d)\varphi(d).\label{eq:quotient}
\end{align}
Thus the finite gcd formulas of A007431, A063659 and A276833
\cite{a007431,a063659,a276833} follow, with local factors
\begin{align*}
 \mathrm{A007431}(p^e)&=
 \begin{cases}p-2,&e=1,\\p^{e-2}(p-1)^2,&e\geq2,\end{cases}\\
 \mathrm{A063659}(p^e)&=
 \begin{cases}p,&e=1,\\p^e-p^{e-2},&e\geq2,\end{cases}
 \qquad
 \mathrm{A276833}(n)=\prod_{p\mid n}(2-p).
\end{align*}
In addition $\sum_{d\mid n}\mathrm{A007431}(d)=\varphi(n)$.
\end{corollary}
\begin{proof}
There are $\varphi(n/d)$ integers $k\leq n$ with $(k,n)=d$,
proving \eqref{eq:gcd}; replacing $d$ by $n/d$ proves
\eqref{eq:quotient}. Taking $F=\mob$ or $F=\mob^2$ gives the first two
sequences. Their prime-power values are respectively
$\varphi(p^e)-\varphi(p^{e-1})$ and
$\varphi(p^e)+\varphi(p^{e-1})$; use $\varphi(1)=1$.
For A276833 take $F=\mob$ in \eqref{eq:quotient} and use
$\operatorname{lcm}(k,n)/k=n/(k,n)$. Factoring its divisor sum yields
$\prod_{p\mid n}(1-\varphi(p))$. All three divisor sums are
multiplicative, since their factors are multiplicative.
Finally, summing $\sum_{d\mid n}\mob(d)\varphi(n/d)$ over divisors
cancels all but $\varphi(n)$.
\end{proof}

\begin{corollary}[Three recorded triangle specializations]\label{cor:triangle-special}
For every $n\geq1$, with $U$ as in Corollary~\ref{cor:triangle},
\[
 a(n)=U(n^2,n)
 =\sum_{\substack{d\mid n\\\mob(d)\ne0}}U(n,d),
 \qquad
 \mathrm{A372706}(n)=\mob(n)U(n^2,n).
\]
These are the three formulas in A400203's history, revision 4,
21 September 2026 \cite{a400203}.
\end{corollary}
\begin{proof}
The first equality uses $\lfloor n^2/n\rfloor=n$ in the definition of $U$.
The divisor sum is \eqref{eq:first} with $X=n$.
The product relation is \eqref{eq:product} with $X=n$.
\end{proof}

\section{Coverage and verification}
Table~\ref{tab:coverage} separates the posed identities, nearby finite
formulas, and analytic questions. The arbitrary-function transform and
rectangular identity are written generalizations of the same cancellation.

\small
\begin{longtable}{@{}p{0.37\linewidth}p{0.57\linewidth}@{}}
\caption{Source statements and mathematical coverage.}\label{tab:coverage}\\
\toprule
Source item & Coverage\\
\midrule
\endfirsthead
\toprule Source item & Coverage\\ \midrule
\endhead
\bottomrule\endfoot
A112399, rev. 41, FORMULA, first conjecture &
Reproved; prior proof Sander--Sander (2009), (5)--(7), specialized;
Ramar\'e (2014), \S1, in equivalent reciprocal-weight form.
Equation~\eqref{eq:first} gives the cutoff form.\\
A112399, rev. 41, FORMULA, second conjecture &
Reproved; prior proof Ramar\'e (electronically 2014; print 2015),
(16.1), reindexed. Coefficient proofs: Ingix/reuns and Helfgott (2019);
later general statements in the March and October 2026 papers.\\
A112399 transform; A332685 transform and lcm formulas &
Reproved; prior proof West (2013), by reduced fractions;
specific formulas recorded in the entries. Corollary~\ref{cor:group}
also states the independent-cutoff form.\\
A372706 product conjecture and nonsquarefree case &
Reproved; prior proof Bala (2026), Proposition~1;
\eqref{eq:product} gives the unweighted formula for arbitrary $X$.\\
A372700, rev. 44, FORMULA, quadratic, row, plateau &
Reproved from the prior finite expansion of Sander--Sander (2009);
the square formula was recorded by Morken (2026), the row/plateau
formulas by Manyama (2024). The rectangular form is \eqref{eq:rectangle}.\\
A372700, rev. 44, COMMENT, infinitely many sign changes &
Not covered.\\
A400203, rev. 17, FORMULA, all ten finite identities &
Proved: Corollary~\ref{cor:triangle}; formulas recorded by Morken
on 21 September 2026, history revisions 2 and 4.\\
A400203, history rev. 4, three additional formulas &
Proved: Corollary~\ref{cor:triangle-special}; recorded by Morken
on 21 September 2026.\\
A112526, COMMENT, signed-variant pointwise product &
Refuted-and-corrected: Corollary~\ref{cor:signed}; the stated local
values require $H(n)\mob(\rad(n))$, not $H(n)\mob(n)$.\\
A073311, FORMULA, A304575 inverse transform &
Reproved; prior general proof West (2013), with $F=\mob^2$;
specific inverse recorded by Oudra (2025).\\
A304575, rev. 19, FORMULA and COMMENT, sums, prime powers, bound &
Reproved: general transform West (2013); specific formulas recorded
by Eldar (2024), Oudra (2025), and the bound by Wiseman (2018).
Equations~\eqref{eq:counts}--\eqref{eq:odd} give the derivations.\\
A304575, rev. 19, NAME, literal strict endpoint &
Refuted-and-corrected: its strict divisor sum is $R(n)-1$;
the inclusive version matches the data.\\
A007431, $\sum\mob((k,n))$, divisor and prime-power formulas &
Reproved; prior general proof Schramm (2008), theorem (i), $m=0$;
specific formulas credited in the entry.\\
A063659, $\sum\mob((k,n))^2$, $\mob^2*\varphi$ and prime-power formulas &
Reproved; prior general proof Schramm (2008), theorem (i), $m=0$;
specific formulas credited in the entry.\\
A276833, divisor/lcm sums and product formula &
Reproved; prior general proof Schramm (2008), after divisor relabeling;
specific formulas credited in the entry.\\
A013929, nonsquarefree interpretation used in A372700/A372706 &
Proved: $\mob(n)=0$ is equivalent to nonsquarefreeness.\\
A002321/A008683/A007947 and other auxiliary entries: unrelated
analytic, Fourier, group-theoretic and enumerative assertions &
Not covered; only their defining conventions and the finite formulas
specified above enter the argument.\\
\end{longtable}
\normalsize

\paragraph{Lean declaration.}
The pinned
\href{https://github.com/the-omega-institute/trureturing/blob/4a11a6467d4ba3682c7630cd22c551cd11074cc6/D5/S3/Arith/MoebiusLogarithmic/MorkenFirstDivisorSum.lean}{\texttt{MorkenFirstDivisorSum.lean}}
contains the declaration
\nolinkurl{D5.S3.Arith.MoebiusLogarithmic.MorkenFirstDivisorSum.result},
whose statement and proof are
\eqref{eq:first} with $X=n>0$.
Its natural-number closed intervals, integer M\"obius values,
squarefree-divisor filter, and inner coprimality condition agree with the
definitions above. At the same revision,
\href{https://github.com/the-omega-institute/trureturing/blob/4a11a6467d4ba3682c7630cd22c551cd11074cc6/D5/S3/Arith/MoebiusLogarithmic/CoprimeConvolution.lean}{\texttt{CoprimeConvolution.lean}}
contains
\nolinkurl{D5.S3.Arith.MoebiusLogarithmic.sum_supported_moebius},
which proves the coefficient identity
\eqref{eq:convolution}. Its predicate \texttt{PrimeSupported} says that
every prime divisor of the index divides the parameter; on positive
indices this is precisely $\rad(q)\mid n$.
The finite summation yielding \eqref{eq:second}, the arbitrary-cutoff
extension in \eqref{eq:first}, and the related identities have the written
proofs given here.

\paragraph{Finite checks.}
The integer-arithmetic checker \texttt{verify.py} compares two
M\"obius implementations and tests the formulas against direct sums for
$1\leq n\leq400$, both cutoff identities and \eqref{eq:product} for
$1\leq n\leq60$, $0\leq X\leq120$, and
\eqref{eq:quadratic}--\eqref{eq:row} for $1\leq N\leq120$.
It also checks \eqref{eq:rectangle} for $0\leq X,Y\leq40$,
the new finite transform and gcd formulas for $1\leq n\leq200$,
and every prime power at most $200$ in \eqref{eq:even}--\eqref{eq:odd}.
The strict-endpoint discrepancy is checked for every $n\leq200$.
All triangle identities are checked through $N=160$ and the
signed-variant correction through $n=200$.
The all-input justification is the written proof.

\paragraph{License.} CC BY 4.0 (Creative Commons Attribution 4.0 International).

\clearpage
\begin{thebibliography}{99}
\small
\setlength{\itemsep}{0pt}
\bibitem{a112399} The OEIS Foundation Inc.,
\emph{A112399: Coprime M\"obius sum}, revision 41,
19 September 2026, FORMULA section; conjectures contributed by
Jamie Morken, 18 September 2026.
\url{https://oeis.org/A112399/internal}.
\bibitem{a008683} The OEIS Foundation Inc.,
\emph{A008683: M\"obius function}.
\url{https://oeis.org/A008683}.
\bibitem{a002321} The OEIS Foundation Inc.,
\emph{A002321: Mertens's function}.
\url{https://oeis.org/A002321}.
\bibitem{a007947} The OEIS Foundation Inc.,
\emph{A007947: Squarefree kernel}.
\url{https://oeis.org/A007947}.
\bibitem{a372706} The OEIS Foundation Inc.,
\emph{A372706: Sum of $\mob(nk)$ for $1\leq k\leq n$}, revision 31,
22 September 2026, FORMULA section.
\url{https://oeis.org/A372706/internal}.
\bibitem{bala} P. Bala,
\emph{Proof of a conjecture in A372706}, 22 September 2026,
Proposition 1.
\url{https://oeis.org/A372706/a372706.pdf}.
\bibitem{a332685} The OEIS Foundation Inc.,
\emph{A332685: Sum of $\mob(k/\gcd(n,k))$ for $1\leq k\leq n$},
revision 15, 15 September 2026, NAME, COMMENT and FORMULA sections.
\url{https://oeis.org/A332685/internal}.
\bibitem{a372700} The OEIS Foundation Inc.,
\emph{A372700: Double sum of $\mob(jk)$}, revision 44,
24 September 2026, FORMULA and COMMENT sections.
\url{https://oeis.org/A372700/internal}.
\bibitem{a073311} The OEIS Foundation Inc.,
\emph{A073311: Squarefree numbers in the reduced residue system},
COMMENT and FORMULA sections.
\url{https://oeis.org/A073311/internal}.
\bibitem{a304575} The OEIS Foundation Inc.,
\emph{A304575}, revision 19, 22 October 2025, NAME, COMMENT,
FORMULA and PROGRAM sections; formulas by Amiram Eldar and Ridouane Oudra.
\url{https://oeis.org/A304575/internal}.
\bibitem{a013928} The OEIS Foundation Inc.,
\emph{A013928: Number of positive squarefree numbers $<n$},
NAME and FORMULA sections.
\url{https://oeis.org/A013928/internal}.
\bibitem{a007431} The OEIS Foundation Inc.,
\emph{A007431: $\sum_{d\mid n}\varphi(d)\mob(n/d)$},
FORMULA section; finite gcd formula by Benoit Cloitre, 14 June 2007;
prime-power formula by Vladeta Jovovic, 25 January 2002.
\url{https://oeis.org/A007431/internal}.
\bibitem{a063659} The OEIS Foundation Inc.,
\emph{A063659: Squarefree gcd count}, FORMULA section;
convolution and gcd formulas by Benoit Cloitre, 8 September 2002 and
14 June 2007; prime-power formula by Vladeta Jovovic, 26 July 2001.
\url{https://oeis.org/A063659/internal}.
\bibitem{a276833} The OEIS Foundation Inc.,
\emph{A276833: $\sum_{d\mid n}\mob(d)\varphi(d)$}, FORMULA section;
product formula by Robert Ferreol, 14 March 2020, and lcm formula
by Peter Bala, 16 January 2024.
\url{https://oeis.org/A276833/internal}.
\bibitem{a400203} The OEIS Foundation Inc.,
\emph{A400203: Triangle of coprime M\"obius sums}, revision 17,
1 October 2026, FORMULA section; entry contributed by Jamie Morken,
21 September 2026.
\url{https://oeis.org/A400203/internal}.
\url{https://oeis.org/history?seq=A400203&start=10}.
\bibitem{triangle-aux} The OEIS Foundation Inc.,
\emph{A112526: Characteristic function for powerful numbers}, NAME,
COMMENT and FORMULA;
\emph{A217038: Number of powerful numbers $<n$}, NAME;
\emph{A062563: $\sum_{k\leq n}\tau(k)\mob(k)$}, NAME and FORMULA.
\url{https://oeis.org/A112526/internal};
\url{https://oeis.org/A217038/internal};
\url{https://oeis.org/A062563/internal}.
\bibitem{sander} J. W. Sander and T. Sander,
\emph{On the Kernel of the Coprime Graph of Integers},
Integers 9 (2009), A43, 569--579; published 19 November 2009.
Proof of Proposition 4, (5)--(7).
\url{https://math.colgate.edu/~integers/j43/j43.pdf}.
\bibitem{ramare2014} O. Ramar\'e,
\emph{Explicit estimates on the summatory functions of the M\"obius function
with coprimality restrictions}, Acta Arithmetica 165 (2014), 1--10,
\S1, after (1.2). \url{https://doi.org/10.4064/aa165-1-1}.
\url{https://ramare-olivier.github.io/Maths/mqdex-8.pdf}.
\bibitem{ramare2015} O. Ramar\'e,
\emph{Explicit estimates on several summatory functions involving the
Moebius function}, Mathematics of Computation 84 (2015), 1359--1387,
\S16, (16.1); electronically published 1 December 2014.
\url{https://doi.org/10.1090/S0025-5718-2014-02914-1}.
\url{https://ramare-olivier.github.io/Maths/mcom2914.pdf}.
\bibitem{mse2019} Ingix and reuns, answers to
\emph{Special sum of Moebius function}, Mathematics Stack Exchange,
13 June 2019; reuns, revision 3, 21:12 UTC.
\url{https://math.stackexchange.com/a/3261283}.
\url{https://math.stackexchange.com/a/3261301}.
\bibitem{helfgott} H. A. Helfgott,
\emph{The ternary Goldbach problem}, manuscript, version listed
16 December 2019, Part I, \S5.3.4, 138--139, (5.73).
\url{https://webusers.imj-prg.fr/~harald.helfgott/anglais/book.html}.
\url{https://webusers.imj-prg.fr/~harald.helfgott/anglais/PartI.pdf}.
\bibitem{rz2026} O. Ramar\'e and S. Zuniga-Alterman,
\emph{On a M\"obius double sum}, arXiv:2603.25961,
v1, 26 March 2026; v2, 1 April; v3, 1 July, Lemma 2.5.
\url{https://arxiv.org/abs/2603.25961v3}.
\bibitem{z2026} S. Zuniga-Alterman,
\emph{A recurrence in logarithmic M\"obius sums with coprimality conditions
via Euler--Maclaurin}, arXiv:2610.17720,
v1, 13 October 2026; v2, 26 October, Lemma 7.1.
\url{https://arxiv.org/abs/2610.17720v2}.
\bibitem{west} B. West, answer to
\emph{Sum of primitive roots is congruent to $\mu(p-1)$ using Moebius inversion?},
Mathematics Stack Exchange, 1 January 2013, 23:07 UTC;
opening reduced-fraction identity and its proof.
\url{https://math.stackexchange.com/a/268866}.
\bibitem{schramm} W. Schramm,
\emph{The Fourier Transform of Functions of the Greatest Common Divisor},
Integers 8 (2008), A50; published 24 November 2008,
theorem (i), (2), and Example 5.
\url{https://math.colgate.edu/~integers/i50/i50.pdf}.
\end{thebibliography}
\end{document}
